\begin{enumerate}
\def\labelenumi{\arabic{enumi}.}
\setcounter{enumi}{4}
\item
  \textbf{Conclusions and Lessons Learned}
\end{enumerate}

\subsection{Conclusions}

The Helike \#213 mission demonstrates that a bioinspired autorotative
recovery system (SRAB) is a viable, simpler and more robust
alternative to conventional parachutes for 1P PocketQubes:

\begin{itemize}
\item The SRAB achieves the LASC regulatory impact velocity with a
  fully passive mechanism — no actuators, no pyrotechnics, no
  software dependency.
\item The computational model (4th-order rigid body + BEM + LEV
  correction) reproduces the measured drop-test behavior with
  sufficient fidelity for design decisions, and the full flight
  simulation (Dédalo launch vehicle + SRAB descent) certifies the
  recovery under real GFS conditions.
\item Compared with a conventional parachute, the SRAB reduces
  descent time by 36\% and total drift by 36\%, improving recovery
  accuracy at the cost of a higher — yet compliant — impact energy.
\item The three-filter review (SCSM compliance, Murphy analysis,
  timeline discipline) proved to be an effective engineering
  framework: it caught the non-parachute notification requirement
  (REC 10.2.1), drove the redundancy decisions in firmware and
  hardware, and kept the project aligned with the competition
  deadlines.
\end{itemize}

\subsection{Lessons Learned}

\begin{itemize}
\item \textbf{LASC 2025: optimization without redundancy is a trap.}
  Every subsystem of Helike has a backup path: SD + LittleFS, LoRa +
  local storage, hardware watchdog, passive recovery that does not
  depend on electronics.
\item \textbf{The parachute is the weakest link.} Replacing it with a
  passive aerodynamic mechanism removes the most failure-prone
  element of the recovery chain.
\item \textbf{Iterative drop testing beats simulation alone.} Each
  test cycle (Asa1 $\rightarrow$ Asa2 $\rightarrow$ Asa3) produced a
  measurable improvement in the model and in the design; the
  simulation-validated iteration loop is the core of the project
  methodology.
\item \textbf{Compliance is a design input, not a formality.} The
  SCSM requirements (kill switches, RBF, autonomy, notification)
  shaped the hardware schematic and the firmware FSM from day one.
\item \textbf{Team management:} the three-filter reviews force
  data-driven decisions with evidence artifacts (videos, logs, bench
  tests), which accelerates the transfer of knowledge from senior
  members to new members — every filter review is a training
  opportunity for the next generation of the team.
\end{itemize}
